% Source: https://artofproblemsolving.com/community/c3h2784571_2021_abmc_team_rounds_58__actonboxborough_math_competition
% Stable identity: aops-topic-2784571-q16
Put $B=(0,0)$ with the tangent line $PB$ as the $x$-axis, $P=(3,0)$ and $O=(0,r)$. Since $PA=AB$, $A$ lies on the perpendicular bisector $x=\dfrac32$ of $PB$, say $A=\left(\dfrac32,y\right)$ with $y>0$.
\begin{asy}
usepackage("amsmath");
size(5cm);
defaultpen(fontsize(10pt));
real s=36/7;real y=sqrt(81/28);real r=s/(2*y);
pair B=(0,0),P=(3,0),O=(0,r),A=(1.5,y),M=(A+B)/2;
draw(circle(O,r));draw((-0.8,0)--(3.6,0));draw(P--A--B);draw(O--M,dashed);draw(O--B,dotted);
draw((1.5,0)--(1.5,y),dotted);
dot(O);dot(M);dot(A);dot(B);dot(P);
label("$O$",O,W);label("$B$",B,SW);label("$P$",P,S);label("$A$",A,NE);label("$M$",M,SE);label("$3$",(1.5,0),S);label("$1$",(O+M)/2,NE);
\end{asy}

$A$ on the circle gives $\dfrac94+(y-r)^2=r^2$, i.e. $AB^2=\dfrac94+y^2=2ry$. Let $s=AB^2$; then $r=\dfrac s{2y}$ and $y^2=s-\dfrac94$. Since $M$ is the midpoint of the chord $AB$, $OM\perp AB$ and
\[1=OM^2=r^2-\dfrac s4=\dfrac{s^2}{4y^2}-\dfrac s4=\dfrac{s^2}{4s-9}-\dfrac s4.\]
Multiplying by $4(4s-9)$: $4s^2-s(4s-9)=4(4s-9)$, so $9s=16s-36$ and $s=\dfrac{36}7$. (Then $y^2=\dfrac{81}{28}>0$, and $P$ is indeed outside the circle.) Hence $m+n=36+7=\boxed{43}$.
